%; whizzy paragraph -pdf xpdf -latex ./whizzypdfptex.sh
%; whizzy-paragraph "^\\\\begin{frame}"
% latex beamer presentation.
% platex, latex-beamer でコンパイルすることを想定。 

%     Tokyo Debian Meeting resources
%     Copyright (C) 2009 Junichi Uekawa
%     Copyright (C) 2009 Nobuhiro Iwamatsu

%     This program is free software; you can redistribute it and/or modify
%     it under the terms of the GNU General Public License as published by
%     the Free Software Foundation; either version 2 of the License, or
%     (at your option) any later version.

%     This program is distributed in the hope that it will be useful,
%     but WITHOUT ANY WARRANTY; without even the implied warreanty of
%     MERCHANTABILITY or FITNESS FOR A PARTICULAR PURPOSE.  See the
%     GNU General Public License for more details.

%     You should have received a copy of the GNU General Public License
%     along with this program; if not, write to the Free Software
%     Foundation, Inc., 51 Franklin St, Fifth Floor, Boston, MA  02110-1301 USA

\documentclass[cjk,dvipdfmx,12pt]{beamer}
\usetheme{Tokyo}
\usepackage{monthlypresentation}

%  preview (shell-command (concat "evince " (replace-regexp-in-string "tex$" "pdf"(buffer-file-name)) "&")) 
%  presentation (shell-command (concat "xpdf -fullscreen " (replace-regexp-in-string "tex$" "pdf"(buffer-file-name)) "&"))
%  presentation (shell-command (concat "evince " (replace-regexp-in-string "tex$" "pdf"(buffer-file-name)) "&"))

%http://www.naney.org/diki/dk/hyperref.html
%日本語EUC系環境の時
\AtBeginDvi{\special{pdf:tounicode EUC-UCS2}}
%シフトJIS系環境の時
%\AtBeginDvi{\special{pdf:tounicode 90ms-RKSJ-UCS2}}

\title{とあるWeb企業での\\Debianシステムの使い方。}
\subtitle{大統一Debian勉強会2013}
\author{株式会社サイバーエージェント\\前田耕平}
\date{2013年6月29日}
%\logo{\includegraphics[width=8cm]{images/openlogo-light.eps}}

\begin{document}

\frame{\titlepage{}}

\begin{frame}
 \frametitle{自己紹介}
  \begin{itemize}
   \item 株式会社サイバーエージェント所属
   \item 主にDC運用関連のソフトウェアの開発やってます
   \item Debian JP Project 2011, 2012年度会長でした
   \item 趣味でPythonでツール開発や、Python関連のDebianパッケージのメンテナンスやってます
   \item 仕事と趣味がほぼ一致
   \item 一児（おまめ\footnote{ネットでの通称}）と一匹（こまめ）の父親です
  \end{itemize}
\end{frame}


\begin{frame}
  \frametitle{今回お話しする内容}
  \begin{itemize}
  \item RHEL系メインの環境でDebianシステム増えてるって？
  \item 社内向けのDebianパッケージの作成、配布方法は？
  \item RPMの配布もDebainでやっちゃうって？
  \end{itemize}
\end{frame}


\begin{frame}
\begin{center}
{\Huge 私からの質問です。}
\end{center}
\end{frame}

\begin{frame}{皆さんの職場はどんな環境？}
 \begin{itemize}
  \item[A.] もちろんDebianメイン
  \item[B.] メインじゃないけどDebianも使ってる
  \item[C.] RHEL系メインでDebian使えれば使いたい
  \item[D.] いや、Unix, GNU/Linuxって何？
 \end{itemize}
\end{frame}

\begin{frame}
  \frametitle{今回お話しする内容}
  \begin{itemize}
  \item \textcolor{red}{RHEL系メインの環境でDebianシステム増えてるって？}
  \item \textcolor{gray}{社内向けのDebianパッケージの作成、配布方法は？}
  \item \textcolor{gray}{RPMの配布もDebainでやっちゃうって？}
  \end{itemize}
\end{frame}


\begin{frame}{サイバーエージェントでは？}
 \begin{itemize}
  \item[1.] CentOSメイン。（今もね）
  \item[2.] 転職時は所属チームのみでUbuntu 10.04 LTSを利用
  \item[3.] 今年2月リリースしたDCはユーザ向け\footnote{私達のチームから見たユーザ=
社内＆グループ会社のエンジニア}のサービスにUbuntu 12.04 LTSを利用\footnote{仮想マシンのホストや、コンシューマ向けサービスは主にCentOS}
  \item[4.] 今月Wheezyも利用可能にした
  \item[5.] 今後リリースするサービスはWheezyを利用予定
 \end{itemize}
\end{frame}

\begin{frame}
\begin{center}
{\Huge Q. どうやって増やしたの？}
\end{center}
\end{frame}

\begin{frame}
\begin{center}
{\Huge A. ワシが使いたいから\\自分で増やした。}
\end{center}
\end{frame}


\begin{frame}
\begin{center}
{\Huge 具体的には？}
\end{center}
\end{frame}

\begin{frame}{OpenStack環境提供期}
 \begin{itemize}
 \item[0.] OpenStack(Cactus, Diablo, Essex) Nova, Glance, Keystone, Horizonの環境としてLucid, Preciseを利用
 \item[1.] アメーバ環境の内部DNS\footnote{PowerDNS＆TonicDNSを使用}用にLucid, Precise利用\footnote{OpenStack用のインスタンス名の自動登録が主目的}
 \end{itemize}
※前提：この時期まではアメーバ向けの環境だけが対象
\end{frame}

\begin{frame}{新DCプロジェクト期}
\begin{itemize}
  \item[2.] Preciseを前提に、必要な機能\footnote{認証,監視,監視用アラートメール,パッケージ管理,ログ管理,サーバ構成管理,NW機器管理,SCM, etc.}を検証＆構築＆開発
  \item[3.] チームメンバーにUbuntuの利用薦める
  \item[4.] 基本公式パッケージを利用
  \item[4.] Sidにも無いソフトウェアや開発したツールは基本全て、Debianパッケージ化
  \item[5.] 作成したパッケージはローカルアーカイブで提供＆公式パッケージ化も目指す
\end{itemize}
※前提：新DCからアメーバだけでなく、他事業部やグループ会社向けの環境も対象
\end{frame}

\begin{frame}{リリース後〜現在}
\begin{itemize}
  \item[6.] チーム内でハンズオンでの勉強会\footnote{Pythonパッケージ作成、Debianパッケージ作成, Git関連など}
  \item[＆] 4,5の繰り返し
 \end{itemize}
\end{frame}

\begin{frame}{身も蓋も無いけど、特別なことは何もない}
\begin{center}
{\Huge 地道にできることをちょっとずつ増やしているだけ}
\end{center}
\end{frame}


\begin{frame}
  \frametitle{今回お話しする内容}
  \begin{itemize}
  \item \textcolor{gray}{RHEL系メインの環境でDebianシステム増えてるって？}
  \item \textcolor{red}{社内向けのDebianパッケージの作成、配布方法は？}
  \item \textcolor{gray}{RPMの配布もDebainでやっちゃうって？}
  \end{itemize}
\end{frame}

\begin{frame}{開発および本番環境}
  \begin{itemize}
    \item (VirtualBox上の)Sidで開発\footnote{毎日スナップショット取ると便利}
    \item PreciseやWheezyは、Sid上でpbuilderを利用
    \item 本番環境はClover\footnote{社内開発のツール。Nova＆cobbler＆DNS＆ラック管理などを提供}で物理及びVMにOSインストール
  \end{itemize}
\end{frame}

\begin{frame}{開発とパッケージメンテナンスの流れ}
  \begin{itemize}
    \item[1.] 使用するミドルウェアもSidのパッケージを利用
    \item[2.] 使用するPythonパッケージもSidのパッケージを利用
    \item[3.] パッケージが無ければ、PyPIにあるものを利用
    \item[4.] PyPIになければ自分でツールやライブラリ開発
    \item[5.] 上記いずれもDebianパッケージ化
  \end{itemize}
\end{frame}

\begin{frame}[containsverbatim]{SidでのDebianパッケージ作成}
Pythonで開発したツールの例\footnote{主にPythonで開発しているので}
  \begin{commandline}
$ cd ~/devel/hoge
$ python setup.py sdist ←(1)
$ cp -i dist/hoge-x.x.tar.gz ~/debpkg/hoge/
$ cd ~/devpkg/hoge
$ tar xvf hoge-x.x.tar.gz
$ cd hoge-x.x
$ dh_make --createorig -s ←(2)
$ rm -rf debian
$ cp -a ~/devel/hoge/debian . ←(3)
$ debuild -us -uc ←(4)
$ cd -
$ sudo pbuilder --build hoge_x.x-x.dsc ←(4)
$ cd /var/cache/pbuilder/result
$ sudo piuparts hoge_x.x-x_amd64.changes ←(5)
$ debsign hoge_x.x-x_amd64.changes ←(6)
$ sudo debi hoge_x.x-x_amd64.changes ←(7)
  \end{commandline}
\end{frame}

\begin{frame}{precise, wheezy向けのパッケージメンテナンス}
（開発 and パッケージ化）したツール

\begin{itemize}
  \item[A) の] ソースパッケージをリビルドする
  \item[B) が] 依存するモジュールがパッケージにない
  \item[C) が] 依存するモジュールがパッケージにはあるが古い
\end{itemize}
\end{frame}

\begin{frame}{パターンAの場合}
\begin{itemize}
  \item pbuilderのbase.tgzイメージを２つ用意
    \begin{itemize}
      \item preciseの場合、precise上で作成したものをSidにコピー
      \item wheezyの場合、Sid上でそのまま作成
    \end{itemize}
  \item 一つはシェルログインでの作業環境用
  \item もう一つは最小構成でのクリーンビルド用
\end{itemize}
\end{frame}

\begin{frame}{パターンB, Cの場合}
\begin{itemize}
\item[1.] 作業環境で依存パッケージ\footnote{依存パッケージがSidの公式パッケージにあればそれ使用を、なければ自分で作っておきます}と最終目的のパッケージをリビルド
\item[2.] 最小構成のbase.tgzから別途コピーしたイメージインストールしておく
\item[3.] 2のbase.tgzを使い、最終目的のパッケージをビルド
\end{itemize}
※Tipsは配布資料を参照。
\end{frame}

\begin{frame}{内部のパッケージアーカイブ}
\begin{itemize}
\item 公式パッケージのアーカイブミラーにはapt-mirror\footnote{\url{http://apt-mirror.github.io/}}
\item カスタムビルドおよび独自パッケージにはreprepro\footnote{\url{http://mirrorer.alioth.debian.org/}} \footnote{python-aptで車輪の再開発やりかけたのは内緒 \url{https://github.com/mkouhei/createapt}}
\end{itemize}
※repreproの細かい手順は配布資料を参照。
\end{frame}

\begin{frame}
  \frametitle{今回お話しする内容}
  \begin{itemize}
  \item \textcolor{gray}{RHEL系メインの環境でDebianシステム増えてるって？}
  \item \textcolor{gray}{社内向けのDebianパッケージの作成、配布方法は？}
  \item \textcolor{red}{RPMの配布もDebainでやっちゃうって？}
  \end{itemize}
\end{frame}

\begin{frame}
{\Huge おまけのRPMリポジトリの作成}
\end{frame}


\begin{frame}
{\Huge いや、CentOSメインなんだからおまけではない}
\end{frame}

\begin{frame}{RPMのリポジトリ}
\begin{itemize}
\item 公式リポジトリミラーはrsyncするだけ
\item カスタムビルドパッケージにはパッケージアップロード用のWeb UIを提供
  \begin{itemize}
    \item Web UIにはFlask
    \item アップロード時のRPMのメタ情報取得にはpython-rpm
    \item リポジトリの作成にはDebianのcreaterepoパッケージ
    \item テストデータ用のRPMはDebian上で作成
  \end{itemize}
\end{itemize}
※Debian上でもRPMは作れるんですよ
\end{frame}

\begin{frame}{まとめというか言いたいこと}
  \begin{itemize}
    \item Debian使いたければ自分で動くしかない
    \item 使うパッケージ管理システムは一つに決めて運用する
      \\もっと良いやり方\footnote{Debianパッケージ管理として。}が合ればぜひ教えてください :)
    \item DebianでだってRPM作れるし、管理できるんですよ
  \end{itemize}
\end{frame}

\begin{frame}
\begin{center}
{\Huge 何かご質問ありますか？}
\end{center}
\end{frame}

\end{document}
